\documentclass[11pt, a4paper]{common/document_template}
\usepackage{common}

\begin{document}

\begin{titlepage}
  \begin{center}
	\Huge{Объединенная межвузовская олимпиада по математике}
  \end{center}
\end{titlepage}  

\chapter{2024}

\begin{exercise}
На доске написаны числа $1, 2, 3, \dots, 27$. За одну операцию Саша может
выбрать два числа на доске, стереть их и записать на доску сумму стёртых чисел. Такими операциями он хочет добиться, чтобы на доске не нашлись одно или несколько чисел, сумма которых равна 27 (сумма одного числа - это само это число). Какое наименьшее количество операций придётся сделать Саше?
\end{exercise}

\end{document}